\chapter{The frontier: Cooper's sporadic operators and the \texorpdfstring{$M_7$}{M7}-polarized family}\label{ch:frontier}

This chapter assembles, with every claim marked, what is currently known about
the third-order operators of Cooper's sporadic sequences and their relation to
the $M_n$-polarized K3 families. It is a status report rather than a theory,
and it is included because a textbook that stops where the textbooks stop
leaves the reader without a picture of what an open problem in this subject
looks like. The detailed account, with proofs and certificates, is the
research paper~\cite{Callens2026paper}; here we give the structure and the
marks.

\section{The template theorem}

\begin{theorem}[\Kstatus; \cite{Callens2026paper}, Thm.~4.1]\label{thm:template}
For every $(a,b,c,d)\in\Q^4$, the order-3 $\theta$-operator $\Lthree$ of the
Cooper--Gorodetsky template
\[
  (n+1)^3s(n+1)=(2n+1)(an^2+an+b)s(n)-n(cn^2+d)s(n-1)
\]
satisfies $\Lthree=P_2\cdot\Sym^2(\Ltwo)$ with
$\Ltwo=P_2\theta^2+P_1\theta+P_0$ explicitly given by polynomials in $z$ over
$\Q[a,b,c,d]$, and $\theta(P_2)=2P_1$. The Almkvist--van~Straten self-adjointness
criterion $W\equiv0$ holds identically in the parameters.
\end{theorem}

The uniformity is the point: the literature treats the sporadic operators one
at a time, and Doran's theorem gives \emph{existence} of a symmetric-square
structure under a geometric hypothesis; here the partner is \emph{exhibited}
from the template parameters with no hypothesis at all.
(Warning~\ref{warn:converse} applies with full force.)

\section{The arithmetic separation}

\begin{theorem}[\Kstatus\ except where marked]\label{thm:separation}
Of the three sporadic sequences, $s_7$ is the only one whose partner operator
$\Ltwo$ has an integral holomorphic solution: the $s_{10}$ and $s_{18}$
partners fail integrality at $n=2$ (values $17/2$ and $45/2$)~\Kstatus; the
$s_7$ partner is the formal square root of $\sum s_7(n)z^n$~\Kstatus\
(\artifact{partner\_eq\_sqrt\_s7}), and its integrality follows from
$4\mid s_7(n)$ for $n\ge1$~\Kstatus\ (\artifact{four\_dvd\_s7}), with no
literature input. The modular interpretation of the $s_7$ partner as a weight-1
form is O'Brien's~\cite{OBrien2016} \Lstatus.
\end{theorem}

\section{The computed monodromy lattice}

\begin{statusbox}{Numerical Claim \Nstatus\ (Stream 2, certificates \texttt{C2\_cooper\_s7\_v5/v6})}
Analytic continuation of $\Ltwo(s_7)$ to $60$ digits, symmetric-squared, and
rational recognition of the monodromy entries (residuals $\sim10^{-59}$),
followed by an exact rational lattice pipeline, gives a joint
monodromy-invariant lattice of rank $3$, even, discriminant $-14$, signature
$(2,1)$, discriminant form $\Z/14$ with $q=1/14$; an explicit
$\GL_3(\Z)$ base change identifies it with $\U\oplus\lat{14}=T_7$. In the
current certificate the stage-2 monodromy matrices are \emph{certified}
(ball arithmetic with rigorous tail bounds, each entry the unique rational
of small denominator in its enclosure); the identification of this lattice
with $T(X)$ of the $s_7$ family remains the framework statement
(Dolgachev, Doran) and is not certified.
\end{statusbox}

\begin{theorem}[\Kstatus; the lattice arithmetic]
Everything in Chapter~\ref{ch:lattices}'s worked examples: $T_N$ is even of
discriminant $-2N$; the glue $e\pm Nf$; the swap is an isometry realizing the
Fricke involution on the period; an explicit embedding witness of
$T_7\oplus(\U\oplus\lat{-14})$ into $\U^3$ with index $14$ and the rank-22
assembly with $E_8(-1)^2$; the modular group $\Gamma_0(N)$ and all
Atkin--Lehner elements inside $\Orth^+(T_N)$ with the weight-2 automorphy
factor; the two singular points of $\Lthree(s_7)$ as the images of the Fricke
fixed points under $z(h)$; and the exact complements of the classes
$e-f$, $(2,-4,1)$, $(14,-14,-5)$ in $T_7$ and $e-f$, $(10,-10,-3)$ in
$T_{10}$.
\end{theorem}

\section{Conjecture and evidence}

\begin{statusbox}{Conjecture (Tier B in the program's language)}
The $s_7$ family is $M_7$-polarized; its very general member has $\rho=19$
and $T\cong\U\oplus\lat{14}$; the rank-2 forms at the CM points are as in
Example~\ref{ex:walls-T7}.
\end{statusbox}

The evidence, each piece with its mark:
\begin{itemize}[leftmargin=2em]
\item The operator is a symmetric square \Kstatus\ --- consistent with Doran,
  not implied by it (Warning~\ref{warn:converse}).
\item The computed monodromy lattice is $T_7$ \Nstatus$\to$\Estatus\ --- consistent
  with Dolgachev's $(M_7)^\perp$.
\item The singular points $\{-1,\tfrac1{27}\}$ are the Fricke fixed points of
  $X_0(7)^+$ \Kstatus\ for the rational map, \Estatus\ (order $40$) for the
  parametrization of the family by $h$ --- this is the one check that is
  \emph{not} forced (Warning~\ref{warn:forced}).
\item The CM points of discriminant $-28,-7,-3$ map to $z=\tfrac1{27},-1,\infty$
  \Nstatus\ (recognition at 120 digits, re-checked at 200).
\end{itemize}

\begin{warning}
What would refute the conjecture: an exact monodromy computation giving a
lattice other than $T_7$; a proof that the $s_7$ family has generic Picard
number $20$ or $\le18$; or a rank-19 polarization other than $M_7$ with the
same period domain (this last is exactly the gap Doran's \S6 leaves open).
Nothing so far points that way, and nothing so far closes the door.
\end{warning}

\section{Two rank-3 lattices, not one}

Remark~\ref{rem:two-lattices} deserves a chapter-level flag here, because it
is the one error the program actually made and caught. The symmetric square of
$\SL_2$ acts on the rank-3 lattice of binary forms of discriminant $-4N^2$;
the period representation $\rho$ acts on $T_N$ of discriminant $-2N$. A
working directive once named them interchangeably. The kernel-checked theorem
\artifact{no\_isometry\_G0N\_TN} is the two-line proof that they are not
isometric --- an isometry preserves the discriminant --- and it exists in the
development only because the conflation had occurred. A textbook lesson: when
two objects carry the same group action, ask what invariant distinguishes them
before giving them the same name.

\section{Open problems}

\begin{enumerate}[leftmargin=2em]
\item Make the monodromy exact (connection formulae for $\Ltwo$), turning
  \Nstatus\ into \Estatus\ throughout.
\item Prove or refute the converse of Doran's theorem for order-3 operators
  with the Cooper template's shape; equivalently, classify rank-19 lattice
  polarizations with period domain $\HH/\Gamma_0(n)^+$.
\item Exhibit explicit Weierstrass models for the two $\rho=20$ members of the
  $s_7$ family with $T=A_2$ and $T=\bigl(\begin{smallmatrix}2&1\\1&4\end{smallmatrix}\bigr)$
  and compute their N\'eron--Severi lattices by Chapter~\ref{ch:elliptic}'s
  methods; the first is $X_3$ (Example~\ref{ex:inose-points}), the second is
  not among Inose's special members and would be new to write down.
\item Formalize Theorem~\ref{thm:tate}'s table and the Shioda--Tate formula so
  that Example~\ref{ex:inose-points} becomes \Kstatus.
\item Formalize Proposition~\ref{prop:index} in general (not just its
  instances), and Nikulin's uniqueness theorem for rank $3$.
\end{enumerate}

\section*{Chapter note}

Every mark in this chapter is transcribed from the program's ledger at the
commit of Chapter~\ref{ch:formal}; the numerical claims are Stream 2's and are
cited by certificate name so the reader can find their provenance. The
author's own contribution to the mathematics is confined to what is marked
\Kstatus; the conjecture is stated as a conjecture.
